% Options for packages loaded elsewhere
\PassOptionsToPackage{unicode}{hyperref}
\PassOptionsToPackage{hyphens}{url}
\documentclass[
  11pt,
]{article}
\usepackage{xcolor}
\usepackage[left=3cm,right=3cm,top=2.5cm,bottom=2.5cm]{geometry}
\usepackage{amsmath,amssymb}
\setcounter{secnumdepth}{5}
\usepackage{iftex}
\ifPDFTeX
  \usepackage[T1]{fontenc}
  \usepackage[utf8]{inputenc}
  \usepackage{textcomp} % provide euro and other symbols
\else % if luatex or xetex
  \usepackage{unicode-math} % this also loads fontspec
  \defaultfontfeatures{Scale=MatchLowercase}
  \defaultfontfeatures[\rmfamily]{Ligatures=TeX,Scale=1}
\fi
\usepackage{lmodern}
\ifPDFTeX\else
  % xetex/luatex font selection
\fi
% Use upquote if available, for straight quotes in verbatim environments
\IfFileExists{upquote.sty}{\usepackage{upquote}}{}
\IfFileExists{microtype.sty}{% use microtype if available
  \usepackage[]{microtype}
  \UseMicrotypeSet[protrusion]{basicmath} % disable protrusion for tt fonts
}{}
\usepackage{setspace}
\makeatletter
\@ifundefined{KOMAClassName}{% if non-KOMA class
  \IfFileExists{parskip.sty}{%
    \usepackage{parskip}
  }{% else
    \setlength{\parindent}{0pt}
    \setlength{\parskip}{6pt plus 2pt minus 1pt}}
}{% if KOMA class
  \KOMAoptions{parskip=half}}
\makeatother
\usepackage{longtable,booktabs,array}
\usepackage{caption}
\captionsetup[table]{skip=6pt}
\newcounter{none} % for unnumbered tables
\usepackage{calc} % for calculating minipage widths
% Correct order of tables after \paragraph or \subparagraph
\usepackage{etoolbox}
\makeatletter
\patchcmd\longtable{\par}{\if@noskipsec\mbox{}\fi\par}{}{}
\makeatother
% Allow footnotes in longtable head/foot
\IfFileExists{footnotehyper.sty}{\usepackage{footnotehyper}}{\usepackage{footnote}}
\makesavenoteenv{longtable}
\usepackage{graphicx}
\makeatletter
\newsavebox\pandoc@box
\newcommand*\pandocbounded[1]{% scales image to fit in text height/width
  \sbox\pandoc@box{#1}%
  \Gscale@div\@tempa{\textheight}{\dimexpr\ht\pandoc@box+\dp\pandoc@box\relax}%
  \Gscale@div\@tempb{\linewidth}{\wd\pandoc@box}%
  \ifdim\@tempb\p@<\@tempa\p@\let\@tempa\@tempb\fi% select the smaller of both
  \ifdim\@tempa\p@<\p@\scalebox{\@tempa}{\usebox\pandoc@box}%
  \else\usebox{\pandoc@box}%
  \fi%
}
% Set default figure placement to htbp
\def\fps@figure{htbp}
\makeatother
\setlength{\emergencystretch}{3em} % prevent overfull lines
\providecommand{\tightlist}{%
  \setlength{\itemsep}{0pt}\setlength{\parskip}{0pt}}
%% tools/stamp.tex
%% Provenance footer for PDFs rendered from this project.
%% Vendored by zzcollab; this file is not project-specific. The
%% footer prints the source document and the time of rendering.
%% The git version is defined here too, and so is available to a
%% project that wants it on the page, but it is left off the
%% default footer: it already names the staged copy in share/ and
%% appears in its manifest row. All values are supplied at render
%% time by tools/stamp-render.R, which writes a small generated
%% file that redefines the macros below. The \providecommand
%% defaults keep a bare LaTeX compile working when the document is
%% built without the wrapper.
\usepackage{fancyhdr}
\usepackage{xcolor}
\providecommand{\stampsource}{\detokenize{(source not recorded)}}
\providecommand{\stamptime}{(time not recorded)}
\providecommand{\stampversion}{\detokenize{(version not recorded)}}
\setlength{\footskip}{40pt}
\newcommand{\stampfootcontent}{%
  \footnotesize\color{gray}%
  \parbox[b]{\textwidth}{%
    \stamptime\hfill\thepage\\[2pt]%
    \ttfamily\raggedright\stampsource}}
\pagestyle{fancy}
\fancyhf{}
\fancyfoot[C]{\stampfootcontent}
\renewcommand{\headrulewidth}{0pt}
\renewcommand{\footrulewidth}{0.4pt}
%% The title page uses the `plain` page style; redefine it so the
%% stamp appears there too.
\fancypagestyle{plain}{%
  \fancyhf{}%
  \fancyfoot[C]{\stampfootcontent}%
  \renewcommand{\headrulewidth}{0pt}%
  \renewcommand{\footrulewidth}{0.4pt}}
\renewcommand{\stampsource}{\textasciitilde{}/\allowbreak{}prj/\allowbreak{}res/\allowbreak{}36-pmsimstats-ng/\allowbreak{}pmsimstats-ng/\allowbreak{}analysis/\allowbreak{}report/\allowbreak{}10-interaction-test-calibration/\allowbreak{}bullets.Rmd}
\renewcommand{\stamptime}{2026-08-15 16:52 PDT}
\renewcommand{\stampversion}{\detokenize{bc0b671-wip}}
\usepackage{amsmath}
\usepackage{amssymb}
\usepackage{booktabs}
\usepackage{longtable}
\usepackage{float}
\usepackage[mathlines]{lineno}
\linenumbers
\usepackage{bookmark}
\IfFileExists{xurl.sty}{\usepackage{xurl}}{} % add URL line breaks if available
\urlstyle{same}
% fallback for those not using the hyperref driver hyperxmp:
\makeatletter
\@ifundefined{xmpquote}{\newcommand{\xmpquote}[1]{#1}}{}
\makeatother
\hypersetup{
  pdftitle={Structured summary: Type I calibration of fixed-effect interaction tests in linear mixed models: a staged diagnosis of working-covariance misspecification and a cluster-robust remedy},
  pdfauthor={pmsimstats team},
  hidelinks,
  pdfcreator={LaTeX via pandoc}}

\title{Structured summary: Type I calibration of fixed-effect interaction tests in linear mixed models: a staged diagnosis of working-covariance misspecification and a cluster-robust remedy}
\author{pmsimstats team}
\date{August 15, 2026}

\begin{document}
\maketitle

{
\setcounter{tocdepth}{2}
\tableofcontents
}
\setstretch{1.1}
\emph{This document is a structured, section-by-section summary of the key points argued in the companion manuscript, ``Type I calibration of fixed-effect interaction tests in linear mixed models: a staged diagnosis of working-covariance misspecification and a cluster-robust remedy.'' It is provided as a supplementary reading aid and does not stand on its own; all notation, derivations, simulation detail, and citations are in the main manuscript.}

\bigskip\noindent

\rule{\textwidth}{0.4pt}\bigskip

\section{Introduction}\label{introduction}

\textbf{The Wald test of a mixed-model interaction is correctly calibrated only when the working covariance is correctly specified.}

\begin{itemize}\setlength{\itemsep}{2pt}\setlength{\parskip}{0pt}
\item The estimand is a fixed-effect interaction coefficient.
\item Inference proceeds by dividing the estimate by a model-based standard error and referring to a $t$ distribution.
\item Correct calibration, exact in balanced cases and first-order otherwise, requires the assumed random-effects structure and residual correlation to match the truth.
\end{itemize}

\textbf{Covariance misspecification renders the model-based standard error inconsistent, motivating the cluster-robust standard error.}

\begin{itemize}\setlength{\itemsep}{2pt}\setlength{\parskip}{0pt}
\item Eicker, Huber, and White established that a misspecified working covariance yields an inconsistent model-based standard error.
\item Liang and Zeger introduced generalized estimating equations (GEE), decoupling point estimation from variance estimation via an empirical sandwich.
\item A substantial literature has since refined the sandwich estimator for the small-sample regime.
\end{itemize}

\textbf{The direction and persistence of miscalibration from a plausible but incorrect working covariance are rarely audited and cannot be answered in the abstract.}

\begin{itemize}\setlength{\itemsep}{2pt}\setlength{\parskip}{0pt}
\item Analysts routinely adopt parametric residual correlations on substantive grounds and trust the resulting standard errors.
\item The magnitude, direction, and sample-size behavior of the resulting miscalibration depend on the specific interaction and covariance mismatch.
\item Answering these questions requires examining a concrete analysis.
\end{itemize}

\textbf{This paper contributes a reusable four-channel diagnostic protocol and a worked finding of asymptotic, conservative miscalibration in a longitudinal interaction test.}

\begin{itemize}\setlength{\itemsep}{2pt}\setlength{\parskip}{0pt}
\item The protocol decomposes the Wald statistic into bias, scale, degrees of freedom, and reference-shape channels, supplemented by a model ladder and a sample-size gradient.
\item The worked example reveals a conservative test whose standard error is over-estimated by roughly ten percent.
\item The conservatism is asymptotic and is removed by a cluster-robust standard error.
\end{itemize}

\textbf{The worked example is drawn from a carryover simulation study and serves solely to make the general phenomenon concrete.}

\begin{itemize}\setlength{\itemsep}{2pt}\setlength{\parskip}{0pt}
\item The example is the biomarker-treatment interaction test in the pmsimstats-ng carryover-sensitivity study.
\item The finding is a property of the mixed-model Wald test, not of the particular simulation.
\item Sections 2 through 5 address the protocol, the example, the staged results, and their generalization.
\end{itemize}

\section{A staged diagnostic protocol}\label{a-staged-diagnostic-protocol}

\textbf{Mis-calibration of a Wald test must enter through exactly one of four separable channels, each diagnosable by simulation.}

\begin{itemize}\setlength{\itemsep}{2pt}\setlength{\parskip}{0pt}
\item The test statistic is $T = \hat\theta / \widehat{\mathrm{SE}}$, referred to a $t_\nu$ critical value.
\item Correct calibration requires the realized rejection rate to equal $\alpha$ under $H_0$.
\item The four channels -- bias, scale, degrees of freedom, and reference shape -- are mutually exhaustive and separable.
\end{itemize}

\subsection{The four-channel decomposition}\label{the-four-channel-decomposition}

\textbf{The decomposition is novel not for introducing new measures but for showing that exactly these four quantities compose into the realized Type I error, making the diagnosis exhaustive and localizing.}

\begin{itemize}\setlength{\itemsep}{2pt}\setlength{\parskip}{0pt}
\item Each channel -- bias, scale, degrees of freedom, shape -- has a classical antecedent: the noncentral $t$ of Johnson and Welch, the level distortions from unequal variance and correlation of Box, the approximate degrees of freedom of Satterthwaite and Welch, and the non-normality studies of Pearson, Box, and Scheff\'e.
\item Together they are exhaustive: a discrepancy absent from all four cannot exist.
\item Because the channels enter the expression in distinct, separately measurable roles, the decomposition localizes the fault rather than merely detecting it.
\end{itemize}

\subsection{The working-covariance model ladder}\label{the-working-covariance-model-ladder}

\textbf{A model ladder localizes the covariance component responsible for standard-error miscalibration by tracking $\kappa$ across a graded sequence of working covariances fitted to the same data.}

\begin{itemize}\setlength{\itemsep}{2pt}\setlength{\parskip}{0pt}
\item Each rung substitutes a different working covariance $V_0$ for the identical simulated data.
\item The rung at which $\kappa$ departs from one identifies the responsible component.
\item Fitting the same data at every rung eliminates Monte Carlo noise between models.
\end{itemize}

\subsection{The sample-size gradient}\label{the-sample-size-gradient}

\textbf{The sample-size gradient distinguishes a finite-sample artifact, which attenuates with $N$, from an asymptotic bias, which persists, and the distinction determines the appropriate remedy.}

\begin{itemize}\setlength{\itemsep}{2pt}\setlength{\parskip}{0pt}
\item A finite-sample artifact diminishes as variance parameters are estimated more precisely; Kenward-Roger corrections address this case.
\item An asymptotic bias is a property of the working covariance mismatch and does not attenuate with sample size.
\item A flat calibration profile across a range of $N$ values identifies the asymptotic case.
\end{itemize}

\section{Worked example: an interaction test under carryover}\label{worked-example-an-interaction-test-under-carryover}

\textbf{The worked example is a biomarker-treatment interaction test from an aggregated N-of-1 carryover simulation study, reported in ADEMP structure.}

\begin{itemize}\setlength{\itemsep}{2pt}\setlength{\parskip}{0pt}
\item Participants are repeatedly crossed between active treatment and control in an N-of-1 design.
\item The scientific question is whether a baseline biomarker moderates the treatment effect.
\item Carryover complicates the analysis and the study compares strategies for accommodating it.
\end{itemize}

\subsection{The analysis model}\label{the-analysis-model}

\textbf{The analysis model is an \texttt{nlme::lme} with a random intercept and corCAR1 residual correlation, and the calibration of its interaction test is the question under examination.}

\begin{itemize}\setlength{\itemsep}{2pt}\setlength{\parskip}{0pt}
\item The model includes a participant random intercept and a continuous-time AR(1) residual correlation.
\item The corCAR1 form was adopted for clinical realism and numerical tractability.
\item Whether it is well calibrated for the fixed-effect interaction test is the central question.
\end{itemize}

\subsection{The null configuration}\label{the-null-configuration}

\textbf{The null configuration collapses the two data-generating process (DGP) architectures to a single process, so the calibration diagnosis need not stratify by architecture.}

\begin{itemize}\setlength{\itemsep}{2pt}\setlength{\parskip}{0pt}
\item Calibration is a null-hypothesis property, evaluated at $c_{bm} = 0$ and $\theta = 0$.
\item The primary cell is the Hybrid design at $N = 70$ with a one-week carryover half-life.
\item At the null, both DGP architectures produce bit-identical data.
\end{itemize}

\textbf{The working covariance is structurally incapable of matching the data-generating covariance, making this a prototypical instance of the misspecification problem.}

\begin{itemize}\setlength{\itemsep}{2pt}\setlength{\parskip}{0pt}
\item The simulated symptom score is a mixture of three autoregressive factors with timepoint-dependent variances plus a participant baseline.
\item A single random intercept and single-parameter corCAR1 residual cannot reproduce this structure.
\item The example is therefore one of plausible, conventional, and inexact working covariance specification.
\end{itemize}

\section{Results}\label{results}

\textbf{Simulation runs are sized to keep the Monte Carlo standard error (MCSE) below 0.005 for each diagnostic component, with all replicate-level outputs retained.}

\begin{itemize}\setlength{\itemsep}{2pt}\setlength{\parskip}{0pt}
\item The decomposition and model ladder use 5000 null replicates; the sample-size gradient uses 1500 per point.
\item The cluster-robust verification uses 3000 replicates.
\item At a Type I error near 0.03, the MCSE is approximately 0.0024 at 5000 replicates.
\end{itemize}

\subsection{The conservatism is standard-error over-estimation}\label{the-conservatism-is-standard-error-over-estimation}

\textbf{The decomposition table excludes point-estimate bias and the degrees-of-freedom channel, leaving the standard-error channel as the site of the conservatism.}

\begin{itemize}\setlength{\itemsep}{2pt}\setlength{\parskip}{0pt}
\item All three specifications show a negligible null-mean displacement of about 0.035 standard deviations.
\item The model assigns approximately 487 degrees of freedom; switching to a normal reference moves the Type I error by only 0.001.
\item The conservatism must therefore reside in the standard-error channel.
\end{itemize}

\textbf{A first-order approximation gives the realized Type I error an elasticity of about $-4.6$ in the standard-error ratio at the nominal point, so a 6--10 percent over-statement is materially consequential.}

\begin{itemize}\setlength{\itemsep}{2pt}\setlength{\parskip}{0pt}
\item With the other three channels inert the realized error reduces to $2\Phi(-\kappa\,t_{\nu,1-\alpha/2})$.
\item The derivative near $\kappa = 1$ is approximately $-0.23$, an elasticity near $-4.6$, so the Type I error falls to the low 0.03 range.
\item The predicted range of 0.031 to 0.037 matches the observed 0.029 to 0.035 to within Monte Carlo error.
\end{itemize}

\textbf{The uniform bowing of the p-value CDF below the diagonal confirms global scale compression rather than a tail anomaly, and the positive-definiteness correction is excluded as a contributor.}

\begin{itemize}\setlength{\itemsep}{2pt}\setlength{\parskip}{0pt}
\item A tail-only failure would depart from the diagonal only near zero; the uniform bowing seen here is distinctive.
\item Global scale compression is the signature of an inflated standard error.
\item The positive-definiteness correction was not triggered in any of the cell's replicates.
\end{itemize}

\subsection{The corCAR1 residual layer is the cause}\label{the-corcar1-residual-layer-is-the-cause}

\textbf{The model ladder identifies the source of the over-estimation by fitting four working covariances to the identical null datasets.}

\begin{itemize}\setlength{\itemsep}{2pt}\setlength{\parskip}{0pt}
\item M0 is ordinary least squares (OLS) with no within-cluster term; M1 adds a random intercept; M2 adds only corCAR1; M3 combines both.
\item Because every rung uses the same data, the comparison is free of Monte Carlo noise.
\item The rung at which $\kappa$ departs from one identifies the responsible component.
\end{itemize}

\textbf{The corCAR1 residual layer is the proximate cause: adding it to a random intercept that is itself well calibrated removes the calibration.}

\begin{itemize}\setlength{\itemsep}{2pt}\setlength{\parskip}{0pt}
\item M1 yields SE ratios of 0.97--0.99 and Type I errors of 0.049--0.055.
\item M3 raises the ratio to 1.07--1.10 and depresses the Type I error to 0.029--0.033.
\item The claim is not that corCAR1 is wrong absolutely but that adding it to a model that already carries a random intercept degrades calibration.
\end{itemize}

\textbf{OLS is strongly conservative for within-cluster contrasts -- reversing the usual Moulton intuition -- and the ladder shows that the sign of the OLS error is a property of the estimand, not of ignoring correlation per se.}

\begin{itemize}\setlength{\itemsep}{2pt}\setlength{\parskip}{0pt}
\item OLS charges the full residual variance to a within-participant contrast when only the within-participant component is relevant, yielding $\kappa \approx 1.3$.
\item The Moulton pitfall applies to between-cluster effects; for within-cluster contrasts the sign is reversed.
\item M2 (corCAR1 alone) is intermediate, confirming that miscalibration magnitude and sign depend on the working covariance choice.
\end{itemize}

\subsection{The conservatism is asymptotic, not a small-sample artifact}\label{the-conservatism-is-asymptotic-not-a-small-sample-artifact}

\textbf{The conservatism is asymptotic: the SE ratio and Type I error show no trend toward their nominal values across an eightfold range of sample size.}

\begin{itemize}\setlength{\itemsep}{2pt}\setlength{\parskip}{0pt}
\item From $N = 35$ to $N = 280$ the SE ratio holds between 1.05 and 1.10 with no downward trend.
\item At $N = 280$ the Type I error is still 0.028, well below the nominal 0.05.
\item The conservatism is a property of the M3 working covariance and will not be outgrown.
\end{itemize}

\textbf{The asymptotic bias arises because the model-based SE converges to a biased target: more data sharpens estimation of the wrong covariance, not the right one.}

\begin{itemize}\setlength{\itemsep}{2pt}\setlength{\parskip}{0pt}
\item The data-generating covariance is a mixture of autoregressive factors with timepoint-dependent variances that M3 cannot reproduce.
\item Larger samples estimate the M3 parameters more precisely but do not correct the structural mismatch.
\item The bias is a property of the model, not of the estimation of the model's parameters.
\end{itemize}

\subsection{A cluster-robust standard error restores calibration}\label{a-cluster-robust-standard-error-restores-calibration}

\textbf{The natural remedy is a cluster-robust SE that is consistent regardless of the working covariance; the CR2 bias-reduced form is chosen for the moderate-cluster regime.}

\begin{itemize}\setlength{\itemsep}{2pt}\setlength{\parskip}{0pt}
\item The corCAR1 model and its point estimate are retained; only the standard error is replaced.
\item The CR2 estimator of Bell and McCaffrey targets the moderate-cluster setting and is implemented in \texttt{clubSandwich}.
\item Satterthwaite degrees of freedom from Pustejovsky and Tipton provide the matching small-sample reference.
\end{itemize}

\textbf{The CR2 estimator uses adjustment matrices $A_i = (I - H_{ii})^{-1/2}$ to achieve exact unbiasedness under the working covariance, and Satterthwaite degrees of freedom are recovered from the eigenvalues of the defining quadratic form.}

\begin{itemize}\setlength{\itemsep}{2pt}\setlength{\parskip}{0pt}
\item CR0 is consistent but downward-biased at moderate cluster counts; the jackknife CR3 over-corrects.
\item CR2 takes the symmetric square root of $(I - H_{ii})$, generalizing the HC2 estimator of MacKinnon and White.
\item For \texttt{lme} objects the construction is applied in the GLS-transformed metric following Pustejovsky and Tipton.
\end{itemize}

\textbf{The CR2 test restores the SE ratio to approximately one and the Type I error to approximately 0.05; Satterthwaite degrees of freedom of approximately 24 confirm the small-sample reference is doing real work.}

\begin{itemize}\setlength{\itemsep}{2pt}\setlength{\parskip}{0pt}
\item CR2 SE ratios are 0.97--1.00 versus 1.05--1.10 for the model-based test.
\item The Type I error rises from approximately 0.03 to 0.045--0.049, within Monte Carlo error of 0.05.
\item Satterthwaite df of approximately 24 (versus model-based 487) indicates that a naive normal reference would be mildly anti-conservative.
\end{itemize}

\textbf{The p-value CDF confirms that CR2 restores calibration without altering the point estimate or the analysis model.}

\begin{itemize}\setlength{\itemsep}{2pt}\setlength{\parskip}{0pt}
\item Model-based p-values bow below the diagonal consistently across all three specifications.
\item Cluster-robust p-values lie on the Uniform(0,1) diagonal.
\item Only the standard error and its reference distribution are replaced; the fitted model is otherwise unchanged.
\end{itemize}

\subsection{The remedy at an alternative cell}\label{the-remedy-at-an-alternative-cell}

\textbf{The comparative ranking of specifications is checked at one alternative cell to determine whether recalibration disturbs their order.}

\begin{itemize}\setlength{\itemsep}{2pt}\setlength{\parskip}{0pt}
\item The alternative cell uses $c_{bm} = 0.45$, the parent study's reference non-null configuration.
\item The question is whether recalibration changes not only absolute power but the ordering of specifications.
\item The Architecture-B Hybrid design at $N = 70$ is retained for comparability with the null diagnosis.
\end{itemize}

\textbf{The recalibration acts as a near-uniform level shift that raises absolute power without disturbing the ranking of specifications.}

\begin{itemize}\setlength{\itemsep}{2pt}\setlength{\parskip}{0pt}
\item At the alternative cell the model-based ratio is 1.12--1.15; CR2 brings it to 0.99 for all three specifications.
\item Power rises by roughly three to four percentage points across specifications.
\item The gaps between specifications are preserved: E2's advantage over E1 and E3 is essentially unchanged.
\end{itemize}

\textbf{A comparative study whose claim is a ranking is robust to recalibration because the correction shifts the level without reordering the contrast.}

\begin{itemize}\setlength{\itemsep}{2pt}\setlength{\parskip}{0pt}
\item The cluster-robust correction acts on the standard error, not on the point estimate.
\item Because the miscalibration is similar across specifications, the shift is approximately uniform.
\item Absolute power figures reported under the model-based test are mildly pessimistic but the ranking is unaffected.
\end{itemize}

\subsection{The result is not about autoregressive correlation}\label{the-result-is-not-about-autoregressive-correlation}

\textbf{The mechanism is general: any parametric working covariance capable of mismatching the truth is capable of miscalibrating the interaction test.}

\begin{itemize}\setlength{\itemsep}{2pt}\setlength{\parskip}{0pt}
\item The corCAR1 term is the component implicated here, but the model ladder shows four different ratios for four different working covariances.
\item Mis-calibration from a wrong working covariance is the Eicker-Huber-White result; the worked example adds quantification and decomposition.
\item The specific contribution is a realistic longitudinal case with two non-obvious features.
\end{itemize}

\subsection{Two features that are not obvious in advance}\label{two-features-that-are-not-obvious-in-advance}

\textbf{For a within-cluster contrast the direction of the standard-error bias can be conservative rather than anti-conservative, reversing the usual sandwich-literature intuition.}

\begin{itemize}\setlength{\itemsep}{2pt}\setlength{\parskip}{0pt}
\item The sandwich literature is typically motivated by anti-conservative standard errors from ignored between-cluster correlation.
\item Here the model-based test is conservative: the corCAR1 term, not the absence of modeling, introduced the bias.
\item An analyst who infers conservatism from having modeled the correlation has the logic backward for a within-cluster contrast.
\end{itemize}

\textbf{Kenward-Roger corrections address finite-sample artifacts and are the wrong instrument here because the bias is asymptotic and does not attenuate with sample size.}

\begin{itemize}\setlength{\itemsep}{2pt}\setlength{\parskip}{0pt}
\item An earlier informal reading attributed the conservatism to moderate sample size; the sample-size gradient refutes this.
\item Kenward-Roger corrections operationalize the Kackar-Harville result for finite-sample variance underestimation.
\item An asymptotic mismatch is not corrected by Kenward-Roger; the cluster-robust SE addresses it directly.
\item Applied empirically through \texttt{mmrm}, Kenward-Roger is inert on the one-parameter autoregressive covariance yet corrects a real finite-sample anti-conservatism on an unstructured one.
\end{itemize}

\subsection{When the problem should be expected}\label{when-the-problem-should-be-expected}

\textbf{The conditions for this problem are met by nearly every applied mixed-model interaction test, and routine accompaniment by a calibration check or cluster-robust SE is warranted.}

\begin{itemize}\setlength{\itemsep}{2pt}\setlength{\parskip}{0pt}
\item The conditions are: a fixed-effect interaction, a parametric working covariance not known to be exact, and a within-cluster contrast.
\item These conditions characterize almost every applied mixed-model interaction analysis.
\item A calibration check or cluster-robust SE removes the question at modest additional cost.
\end{itemize}

\subsection{The marginal-model alternative}\label{the-marginal-model-alternative}

\textbf{GEE is a natural alternative that carries the sandwich variance from the outset; for a linear model with identity link it estimates the same interaction as the conditional mixed model.}

\begin{itemize}\setlength{\itemsep}{2pt}\setlength{\parskip}{0pt}
\item Liang-Zeger GEE decouples point estimation from variance estimation by design.
\item For a linear identity-link model, marginal and conditional fixed-effect coefficients coincide.
\item The working correlation governs efficiency, not inferential validity.
\end{itemize}

\textbf{The GEE route is viable but requires its own small-sample correction; both the conditional CR2 and marginal Mancl-DeRouen tests are defensible calibrated analyses and the choice is practical rather than principled.}

\begin{itemize}\setlength{\itemsep}{2pt}\setlength{\parskip}{0pt}
\item The naive GEE sandwich under-estimates the sampling SD by about four percent, yielding a Type I error of 0.067--0.069.
\item Mancl-DeRouen bias correction brings the ratio to within one percent of one and the Type I error to 0.055--0.058.
\item The conditional CR2 test is slightly conservative; the marginal Mancl-DeRouen test slightly anti-conservative; both are defensible.
\end{itemize}

\subsection{Limitations}\label{limitations}

\textbf{The scope is bounded to one data-generating family, one estimand, and one working covariance, and the power consequence is verified at only a single alternative cell.}

\begin{itemize}\setlength{\itemsep}{2pt}\setlength{\parskip}{0pt}
\item Generality beyond the corCAR1-plus-random-intercept model rests on the model ladder and established theory, not on an empirical sweep.
\item The gradient varied the number of clusters but not the cluster size; a gradient in timepoints per participant is a further extension.
\item The CR2 remedy spends precision (Satterthwaite df approximately 24) to buy validity, and that cost rises at smaller cluster counts.
\item Whether the near-uniform power shift holds across the full factorial rather than the one cell examined remains an open question.
\end{itemize}

\section{Reproducibility}\label{reproducibility}

\textbf{All diagnostic outputs are generated by seven named scripts and read inline; the manuscript is fully reproducible from the repository root.}

\begin{itemize}\setlength{\itemsep}{2pt}\setlength{\parskip}{0pt}
\item Seven scripts under \texttt{analysis/scripts/carryover-sensitivity/} cover the decomposition, model ladder, sample-size gradient, null and alternative CR2 verifications, the GEE check, and the Kenward-Roger check.
\item Each script writes a replicate-level \texttt{.rds} file; this manuscript reads those files and regenerates every table and figure inline.
\item The examination plan that preceded the manuscript is retained alongside it as \texttt{type1-examination-plan.md}.
\end{itemize}

\end{document}
